\documentclass{article}
\usepackage[T1]{fontenc}
\usepackage{lmodern}
\usepackage{graphicx}
\usepackage{amsmath,amssymb}
\usepackage[a4paper,margin=2cm]{geometry}
\usepackage[absolute,overlay]{textpos}
\setlength{\TPHorizModule}{1mm}
\setlength{\TPVertModule}{1mm}
\usepackage[indonesian]{babel}
\usepackage{hyperref}
\setlength{\emergencystretch}{2em}
% unit: Halaman 16

\begin{document}

% === Halaman 16 (python/native) ===

• Sebagai latihan, misalkan diberikan sebuah cipherteks:


TLCYKUMGDFAWTZVOYKLENSZZHYZRW


    temukan kunci yang menghasilkan plainteks:

       mr johnson left his house last night

     lalu temukan kunci lain yang menghasilkan  plainteks 


i saw the mysterious plane behind me

16

\end{document}
